\scanpage{42}
ment'' évidente, qui ici dit \deleted{\unsure{qu'}}\inserted{que la donnée}
\unsure{d'une} gerbe \inserted{$\unsure{\mathcal G|G}_X$}
sur $X$, revient à celle de gerbes $\unsure{\mathcal G|G}_{X^*}$ sur
$X^*$ et $\unsure{\mathcal G|G}_{Y,X}$ sur $V_{Y,X}$, et d'une
\deleted{\unsure{iso.}}\inserted{équiv.} de \deleted{\unsure{recollem.}} gerbes
\[
(24)\qquad \unsure{\mathcal G|G}_{Y,X}\mid V^*_{Y,X}\simeq
\unsure{\mathcal G|G}_{X^*}\mid V^*_{Y,X}.
\]
Or sur un $K(\pi,1)$, une $A$-gerbe
s'interprète canoniquement comme une
\uline{extension} de $\pi$ par $A$. Donc la donnée
d'une gerbe sur $X\deleted{{}\simeq S^2}$ revient à celle
d'\deleted{\unsure{une ext}} deux extensions, l'une de
$\pi_1(Y)$ par $A$, l'autre de $\pi_1(X^*)$ par $A$
et d'un isom. entre les deux extensions
de $D$ par $A$ qui s'en déduisent
par image inverse par les hom. (10).
\margin{\diagram{\begin{tikzcd}[column sep=small,row sep=small]
D\arrow[r]\arrow[d]&\pi_1(Y)\\
\pi_1(X^*)&
\end{tikzcd}}}
(cette description marche dès que les
``pièces'' $Y$, $X^*$, $V^*_{Y,X}$ sont connexes, et
que leur $\pi_2$ est nul). Dans le cas
\mspage{41}
